\documentclass[10pt]{article}
\usepackage{titling}

\title{Conformational Analysis}
\date{DATE PLACEHOLDER}
\author{Ingyu Bahng}

\usepackage[letterpaper, margin=1in]{geometry}
\usepackage{amsmath} % better math functions
\usepackage{amssymb} % better math symbols
\usepackage{babel}[english] % change rules to english
\usepackage{lipsum} % allows for lorem ipsum text generation
\usepackage{xr-hyper} % allows cross referencing labels from external LaTeX documents 
\usepackage{hyperref} % for hyperlinking text
\usepackage{fancyvrb} % for better verbatim environments
\usepackage{newverbs} % allows for inline typewriter font via \texttt{}
\usepackage{fancyvrb} % also code font blocks but has more capabilities than texttt
\usepackage{footnote} % allows footnotes inside tabulars
\usepackage{dcolumn} % allows for decimal alignment in tables
\usepackage{tabularx}

\usepackage{fontspec}
\setmainfont{Gentium Plus} % allowing greek letters in text

\usepackage{unicode-math}
\setmathfont{XITS Math} % allowing greek letters in math



% tcolorbox package
\usepackage{tcolorbox}
\tcbuselibrary{breakable}

\tcbset{ % this is the default design of tcolorboxes
  colframe = black,
  colback  = black!1,
  coltitle = black,
  colbacktitle = black!15,
  breakable, 
  arc=0mm,
  boxrule=0.5pt,
  toptitle=3pt,
  bottomtitle=3pt,
  left=8pt,
  right=8pt,
  top=6pt,
  bottom=6pt,
  before skip=12pt,
  after skip=12pt,
  bottomrule at break=-1pt,
  toprule at break=-1pt,
  fonttitle=\bfseries,
}

% Custom Environments
\newtcolorbox[auto counter, number within=section]{definition}[1][]
{
    title = \textbf{Definition \thetcbcounter: ~#1}
}

\newtcolorbox[auto counter, number within=section]{core}[1][]
{
    title = \textbf{Core Concept \thetcbcounter: ~#1}
}


\usepackage{fancyhdr}
\usepackage[skip=0.8em, indent=0em]{parskip}

\pagestyle{fancy} % this section generates the lines and text at the top and bottom of each page
\fancyhead[L]{\thetitle}
\fancyhead[C]{\theauthor}
\fancyhead[R]{\thedate}
\fancyfoot[C]{\thepage}
\renewcommand{\footrulewidth}{0.3pt}
\renewcommand{\thispagestyle}[1]{}

\renewcommand{\arraystretch}{1.5} % table vertical padding
\setlength{\tabcolsep}{10pt} % table horizontal padding

\usepackage{enumitem} % better control over enumerate and itemize

% List customization
\setlist[itemize]{%
  label=\bullet,
  itemsep=2pt,
  leftmargin=1.5em,
}

\setlist[enumerate]{%
  label=(\alph*),
  itemsep=2pt,
  leftmargin=1.5em,
}



\usepackage{pgfplots} % for plot creation
\pgfplotsset{compat=1.18} % setting the version used for pgfplots
\usepackage{float} % for better figure placement using [h]
\usepackage{graphicx} % for importing figures into the tex file
\usepackage{subcaption} % allows for multiple subfigures with individual captions wihtin one figure
\usepackage{xparse} % allows for more than 9 parameters in commands
\usepackage{adjustbox} % for valign option
\captionsetup{font=small} % setting the caption fontsize to small always

\usepackage{silence}
\WarningsOff[float]  % suppress float package warnings

\newcommand{\myfig}[4]{%
  \begin{figure}[H]
    \centering
    \adjustbox{valign=c}{\includegraphics[width=#1\textwidth]{#2}}
    \caption{#3}
    \label{#4}
  \end{figure}
}

\newcommand{\mymultifig}[8]{
  \begin{figure}[H]%
    \centering%
    \begin{subfigure}{#1\textwidth}%
      \centering%
      \includegraphics[width=\linewidth]{#2}%
    \end{subfigure}%
    \hfill%
    \begin{subfigure}{#3\textwidth}%
      \centering%
      \includegraphics[width=\linewidth]{#4}%
    \end{subfigure}%
    \hfill%
    \begin{subfigure}{#5\textwidth}%
      \centering%
      \includegraphics[width=\linewidth]{#6}%
    \end{subfigure}%
    \caption{#7}%
    \label{#8}%
  \end{figure}%
}

\usepackage{newverbs} % allows for inline typewriter font via \texttt{}
\usepackage{fancyvrb} % also code font blocks but has more capabilities than texttt
\usepackage{listings} % allows for codeblocks - config below

%BELOW ARE CODE BLOCK DESIGNS (Light Theme)
\definecolor{gray}{HTML}{707070}
\definecolor{lightgray}{HTML}{333333}
\definecolor{codebg}{HTML}{F5F5F5}
\definecolor{keyword}{HTML}{0000FF}
\definecolor{string}{HTML}{A31515}
\definecolor{number}{HTML}{098658}
\definecolor{function}{HTML}{795E26}
\definecolor{comment}{HTML}{008000}
\definecolor{classname}{HTML}{267F99}

% default options for listings (for code)
\lstset{
  frame=ltbr,
  language=python,
  aboveskip=3mm,
  belowskip=3mm,
  showstringspaces=false,
  columns=fullflexible,
  keepspaces=true,
  basicstyle=\fontsize{9}{10}\ttfamily\color{lightgray},
  numbers=left,
  firstnumber=1,
  numberstyle=\fontsize{7}{10}\ttfamily\color{gray},
  keywordstyle=\color{keyword},
  commentstyle=\color{comment},
  stringstyle=\color{string},
  backgroundcolor=\color{codebg},
  breaklines=true,
  breakatwhitespace=true,
  tabsize=2,
  xleftmargin=3em,
  numbersep=1em,
  framexleftmargin=2.5em,
  framesep=6pt,
  stepnumber=1,
}

\hfuzz=5.002pt % ignore overfull hbox badness warnings below this limit



\usepackage{chemfig}
\usepackage[version=4]{mhchem}

\newcommand{\bce}[1]{
  \hspace{7mm}{#1}
}

\newcommand{\sno}[0]{
  $\mathrm{S}_{\mathrm{N}}1$
}

\newcommand{\snt}[0]{
  $\mathrm{S}_{\mathrm{N}}2$
}

\newcommand{\reactionfig}[4]{
    \begin{figure}[H]
    \centering
    \scalebox{#4}{
    \schemestart
    #1
    \schemestop
    }
    \caption{#2}
    \label{#3}
    \end{figure}
}




\begin{document}

\input{../../presets/_general/startup}

\subsection{Rotamers}

  \textbf{Rotamers}, short for rotational isomers, are a specific type of conformational isomer. They are different spatial arrangements of the same molecule that arise purely from rotation around a single ($\sigma$) bond.

  As we have repeatedly seen throughout the \textit{Introductory Chemistry} notes, matter arranges itself in a way that minimizes its PE. Molecule arrangements are no different.

  Let's take a look at different conformational isomers of ethane through a \textbf{Newman projection}.

  \myfig{0.9}{conformational_analysis_figures/ethane_rotamer1}{caption}{fig:fig1}

  As shown above, there are two different types of positions ethane can take when rotated about its C-C $\sigma$ bond: \textbf{eclipsed} and \textbf{staggered}. Of the two, the most stable rotamer is the staggered one because the \ce{e^-} \ce{e^-} repulsions between the C-H $\sigma$ bonds are minimized.

  On the flipside, the eclipsed conformation is the least stable because the \ce{e^-} \ce{e^-} repulsions between the C-H $\sigma$ bonds are maximized.

  The PE diagram throughout a full rotation around the C-C $\sigma$ bond looks like the following.

  \myfig{0.7}{conformational_analysis_figures/ethane_rotamer2}{caption}{fig:fig2}

  Here, the $3.0kcal\,mol^{-1}$ would represent the activation energy ($E_a$) to carry out a rotation from a staggered to an eclipsed conformation. Since the different staggered positions of ethane are exactly identical, there is no net change in the energy of the system after a rotation occurs (neither endothermic nor exothermic). We also see the same pattern in the case of propane.

  \myfig{0.7}{conformational_analysis_figures/propane_rotamer1}{caption}{fig:fig3}

  The only difference we see is the $E_a$ increasing to $3.4\,kcal\,mol^{-1}$ for propane. This is because there is greater \ce{e^-} density within the \textbf{methyl of propane} compared to the \textbf{hydrogen of ethane}. Thus the \ce{e^-} \ce{e^-} repulsion forces are greater.

  Let's take a look at one more example: butane.

  \myfig{0.75}{conformational_analysis_figures/butane_rotamer1}{caption}{fig:fig4}

  Using the same logic we learned from the greater \ce{e^-} density consequences in propane compared to ethane, we can see that the anti staggered conformation would be the most stable becasue the methyl groups are placed farthest apart. Second most stable would be the gauche staggered conformation. The fully eclipsed conformation ($\theta = 0^{\circ}$) would be the least stable because the high \ce{e^-} density methyls are placed exactly overlapping on one another, whereas the half eclipsed conformation ($\theta = 120^{\circ}$) would be more stable but still unstable in the grand scheme of the molecule's PE levels.

  \myfig{0.75}{conformational_analysis_figures/butane_rotamer2}{caption}{fig:fig5}

  \begin{definition}
  \textbf{Stereoisomers} are different spatial arrangements of the same molecule (same connectivity) that arise from differences in 3D arrangements. Two types of stereoisomers exist: \textbf{conformational} and \textbf{configurational} isomers. \\

  \textbf{Conformational isomers} are the result of bond rotations that happen freely at room temperature. \textit{Rotamers are conformational isomers.} \\

  \textbf{Configurational isomers} are permanent, meaning they can only be change by the breaking and formation of new bonds. They refer to the different spatial arrangements like cis/trans and R/S.
  \end{definition}

\subsection{Cycloalkane \& Ring Strain}

  Cycloalkanes experience ring strain (\textbf{heightened PE}) due to the following factors.

  (1) \textbf{Bond angle strain} (\textit{especially for smaller rings})

  \myfig{0.5}{conformational_analysis_figures/cycloalkane_bond_angles}{caption}{fig:fig6}

  (2) \textbf{Eclipsing strain} (\textit{as we have discussed in 1.2.1 Rotamers})

  (3) \textbf{Transannular strain} (\textit{the steric hinderence that occurs from atoms "bumping into each other" across the empty space in the middle of the ring})

  \myfig{0.35}{conformational_analysis_figures/transannular_strain}{caption}{fig:fig7}

  Thus, we see a variety of 3D arrangements across cycloalkanes that differ from the flat ring orientation that we intuitively expect.

  \myfig{0.55}{conformational_analysis_figures/ring_arrangements}{caption}{fig:fig8}

  Specifically, we will be focusing on the different conformational arrangements of \textbf{cyclohexane}, which most stably exists in a chair shape conformation. 

  \myfig{0.5}{conformational_analysis_figures/cyclohexane_chair}{caption}{fig:fig9}

  \textbf{MINIMIZING ECLIPSING STRAIN:} One reason why this conformation is the most stable is because the of the staggered rotamer conformation when viewing the Newman projection along any of the C-C bonds as we have learned in \textbf{1.2.1 Rotamers}, eliminating the eclipsing strain experienced \textit{if} the cyclohexane took a planar conformation.

  \myfig{0.65}{conformational_analysis_figures/cyclohexane_newman}{caption}{fig:fig10}

  \textbf{MINIMIZING TRANSANNULAR STRAIN:} Another reason why the chair conformation is favored is because the conformation minimizes steric hinderence across the ring. \textit{Eclipsing strain is also increased in this conformation.} Take a look at the \textbf{boat conformation} of cyclohexane for instance.

  \myfig{0.5}{conformational_analysis_figures/cyclohexane_transannular}{caption}{fig:fig11}

  The boat arrangement is actually a transition state in the dynamics of cyclohexane movement. Thus it does not exist readily in the environment. Instead, a slightly adjusted, more stable version is the \textbf{twist (skew) boat conformation}.

  \myfig{0.6}{conformational_analysis_figures/twist_boat_transition}{caption}{fig:fig12}

  Here is a PE graph across different conformations of cyclohexane.

  \myfig{0.8}{conformational_analysis_figures/cyclohexane_pe}{caption}{fig:fig13}

  As you can see, there is another high energy (\textit{for the same reasons we have discussed}) transition state, the \textbf{half chair conformation}.

  While nonsubstituted cyclohexane shows indifference to the two different chair conformations in regards to PE levels, it's important to keep in mind that monosubstituted cyclohexanes (such as methylcyclohexane) do show a preference between the two chair conformation due to differing levels of transannular strain.

  \myfig{0.65}{conformational_analysis_figures/methylcyclohexane}{caption}{fig:fig14}

  In other words, the \textbf{equitorial} methyl chair conformation is preferred over the \textbf{axial} methyl chair conformation due to the lower transannular strain.

  \begin{definition}

  \textbf{Axial} bonds refer to bonds that point straight up or straight down, perpendicular to the "plane" of the ring. On the other hand, \textbf{Equatorial} bonds are those that point outward toward the "equator" or perimeter of the ring.

  \myfig{0.35}{conformational_analysis_figures/eq_ax}{caption}{fig:fig15}

  \end{definition}

\end{document}
